




\documentclass[11pt]{article}

\usepackage[a4paper,margin=1in]{geometry}
\usepackage[T1]{fontenc}
\usepackage[utf8]{inputenc}
\usepackage{lmodern}
\usepackage{amsmath,amssymb,amsthm,bm,mathtools}
\usepackage{graphicx}
\usepackage{float}
\usepackage{placeins}
\usepackage{enumitem}
\usepackage{array}
\usepackage{longtable}
\usepackage{tabularx}
\usepackage{booktabs}
\usepackage{xurl}
\usepackage[expansion=false]{microtype}
\setlength{\emergencystretch}{3em}
\usepackage{xcolor}
\usepackage{hyperref}
\usepackage[nameinlink,noabbrev]{cleveref}

\usepackage[backend=biber,style=numeric,sorting=none,maxnames=6,uniquename=false,url=false]{biblatex}
\AtEveryBibitem{\clearfield{note}\clearfield{annotation}\clearfield{keywords}}
\addbibresource{cmotion_foundational_general.bib}
\addbibresource{cmotion_primary_research.bib}
\addbibresource{cmotion_internal_manuscripts.bib}

\newtheorem{theorem}{Theorem}[section]
\newtheorem{lemma}[theorem]{Lemma}
\newtheorem{proposition}[theorem]{Proposition}
\newtheorem{corollary}[theorem]{Corollary}
\theoremstyle{definition}
\newtheorem{definition}[theorem]{Definition}
\newtheorem{axiom}[theorem]{Foundational commitment}
\newtheorem{remark}[theorem]{Remark}

\newcommand{\cmotion}{\mbox{C-Motion}}
\newcommand{\ElEmEr}{\mbox{ElEmEr}}
\newcommand{\cmc}{c}
\newcommand{\dd}{\mathrm{d}}
\newcommand{\ii}{\mathrm{i}}
\newcommand{\RR}{\mathbb{R}}
\newcommand{\CC}{\mathbb{C}}
\newcommand{\UU}{\mathrm{U}(1)}
\newcommand{\SU}{\mathrm{SU}(2)}
\newcommand{\SO}{\mathrm{SO}(3)}
\newcommand{\Order}{\mathcal{O}}

\newcommand{\statusDerived}{\textcolor{green!50!black}{\textbf{Derived}}}
\newcommand{\statusDerivedConst}{\textcolor{green!35!black}{\textbf{Derived*}}}
\newcommand{\statusModelDerived}{\textcolor{cyan!55!black}{\textbf{Model-derived}}}
\newcommand{\statusStructural}{\textcolor{blue!55!black}{\textbf{Structural}}}
\newcommand{\statusNumerical}{\textcolor{blue!60!black}{\textbf{Numerical}}}
\newcommand{\statusMatched}{\textcolor{orange!85!black}{\textbf{Matched}}}
\newcommand{\statusAssumed}{\textcolor{orange!80!black}{\textbf{Assumed}}}
\newcommand{\statusConditional}{\textcolor{yellow!60!black}{\textbf{Conditional}}}
\newcommand{\statusNegative}{\textcolor{red!65!black}{\textbf{Negative}}}
\newcommand{\statusConjectural}{\textcolor{magenta!65!black}{\textbf{Conjectural}}}
\newcommand{\statusDeferred}{\textcolor{gray!70!black}{\textbf{Deferred}}}
\newcommand{\statusOpen}{\textcolor{red!75!black}{\textbf{Open}}}
\newcommand{\statusSharpOpen}{\textcolor{red!35!black}{\textbf{Open (sharp)}}}
\newcommand{\statusRetracted}{\textcolor{red!80!black}{\textbf{Retracted}}}

\newcolumntype{L}{>{\raggedright\arraybackslash}X}
\newcolumntype{P}[1]{>{\raggedright\arraybackslash}p{#1}}


\hypersetup{
  colorlinks=true,
  linkcolor=blue!50!black,
  citecolor=blue!50!black,
  urlcolor=blue!50!black,
  pdftitle={C-Motion: Local-c Propagation, Recurrent Closure, and Field-First Interaction},
  pdfauthor={Toma Kavonius}
}

\title{\textbf{C-Motion: Local-\texorpdfstring{$c$}{c} Propagation, Recurrent Closure, and Field-First Interaction}}
\author{Toma Kavonius}
\date{24 November 2026}


\begin{document}
\maketitle

\begin{center}
\small\itshape
Scope and limitations.  This manuscript presents the current bounded state of the
C-Motion framework.  Open assumptions, conditional results, and unresolved
completion conditions are stated explicitly and are not presented as completed
microscopic derivations.
\end{center}
\vspace{0.6em}

\begin{abstract}
C-Motion investigates whether locally universal-$c$ propagation in physical
three-space can support physical clocks, persistent matter-like organization,
and interaction without treating clock time as a primitive material dimension.
Its proposed underlying field/state structure, Elemental Motion Resonance
(ElEmEr), admits open propagating modes and candidate recurrent
\emph{closures} as alternative organizations, not distinct primitive
substances. An assumed succession order $s$ distinguishes physical
before-and-after from operational time measured by recurrent processes and
signals. A moving closure could alter its cycle rate by reorganizing its
return geometry, but the recovery of operational Lorentz kinematics remains
conditional on complete clock, rod, and signal dynamics with no observable
$s$-dependence.

The framework treats interaction as reorganization of the complete local
field state. Its candidate source--response--propagation architecture
separates an interaction opportunity from actual activation, steering,
capture, and sustained binding. Conditional speed-preserving steering and
complete-stress relations constrain possible dynamics; they do not supply a
selected microscopic action. For smooth recurrent symmetry families with an
exact action charge, the energy--charge derivative obeys
$\mathrm{d}E/\mathrm{d}Q=J_0\Omega$ under stated assumptions, while a
generic product law and the identification $J_0=\hbar$ do not follow.
The central research requirements remain a physically selected local law,
finite stable closures, absolute-scale selection, and the empirical particle
spectrum. This paper presents the framework's positive construction,
checkable conditional constraints, and explicit completion tests as ongoing
work rather than a completed microscopic theory.
\end{abstract}

\medskip

\section{Introduction: from local propagation to organized physics}
\label{sec:introduction}

Modern physics describes its best-tested domains with several powerful but conceptually
different structures: spacetime geometry in gravitation, quantum states and measurement
rules in quantum theory, and gauge fields and symmetry breaking in the Standard Model.
One way to seek a deeper account is to add further structure.  \cmotion{} explores the
opposite possibility: how far can one get from a simpler underlying dynamics of
propagation and organization?

The question can be stated physically:
\begin{quote}
\emph{How much familiar physics can arise if primitive physical change consists of
local-$c$ field propagation in physical three-space, while clocks, massive persistence,
and interactions arise from organization of that propagation rather than from treating
time as a fundamental spatial-like dimension?}
\end{quote}

\cmotion{} began from a small set of proposed physical principles and questions about
motion, time, interaction, and recurrent structure.  Much of the present framework is
the result of testing how far those starting ideas can be carried, where they require
refinement, and where additional physical structure becomes unavoidable.

The proposed underlying structure is \ElEmEr{}.  It is not introduced as a collection of
little persistent particles moving through an otherwise passive background.  It is the
field/state content itself.  The same underlying structure may propagate openly, bend or
change mode in a local interaction, return recurrently, or participate in an extended
organization.  A signal, a clock, and a particle-like persistent state would therefore
be different dynamical organizations of one field rather than different primitive kinds
of matter.

Taken together, these commitments make bending and reorganization more basic than
microscopic speeding-up or slowing-down while the primitive propagation magnitude remains fixed.  Recurrence turns ongoing propagation into a physical cycle and hence into
a possible clock.  Sufficiently stable recurrence offers a route from motion to
localized persistence.  Because the local interaction object is the complete field
state, forces and conversions can be sought as local reorganization followed by causal
field extension rather than as instantaneous action between permanently distinct
parcels.

\paragraph{A physical route through the proposal.}
Consider first an outward-moving disturbance.  In the proposed ontology it is an open
organization of the field: it can transport a change without being a miniature
persistent object.  If that disturbance enters a region where the complete field
state activates a response, its direction or internal mode may change while its
resolved primitive propagation remains locally at $c$.  A response can continue
to send an altered disturbance outward, or it could participate in a sequence
of returns.  Nothing in this description assumes that every encounter produces
attraction or that every bent trajectory becomes a closure.

Now imagine that a particular field pattern does repeatedly return in a way that
reconstructs its internal organization, sheds or balances its excess flux, and
remains localized.  Such a \emph{hypothetical dynamically sustained} return would
be the candidate material object.  Its repeated cycle can also be read as a clock;
its whole-pattern energy and momentum, if defined by the native dynamics, are the
quantities relevant to mass and inertia.  Relative motion or an external interaction
could change the geometry of this recurrence and therefore the rate of the clock,
without requiring a primitive field component to travel slower than $c$.

This is a connected physical proposal, not a demonstrated history of closure formation.
The same missing microscopic law must decide which interactions activate, whether
recurrent solutions exist, how energy and stress are carried, and which organizations
survive.  The mathematical and empirical constraints in the following sections
specify what that proposed law would have to accomplish; they do not replace the
physical picture it is meant to describe.

The construction is developed in that order.  We first state the physical commitments,
then follow their immediate consequences for clocks, interaction, steering, and closure.
Energy, action, stress, operational relativity, and scale constraints are introduced only
when the preceding structure requires them.  Quantum normalization, thermodynamics, gravity,
gauge structure, and entanglement are treated in companion papers where their additional
assumptions and calculations can be stated explicitly. Electroweak mode selection and
normalization remain subjects of ongoing research.

\section{Foundational physical commitments}
\label{sec:premises}

The following commitments are the present distilled foundation of \cmotion{}.  Many trace
back to the framework's earlier physical proposals, although their meaning and scope have
been sharpened by the subsequent derivations and tests.  The next section asks what
becomes possible once they are taken together; the formal claim classifications are
collected in \cref{app:status}.

\subsection{Elemental Motion Resonance in physical three-space}
\label{subsec:three_space}

\cmotion{} takes physical three-space as primitive.  \textbf{Elemental Motion Resonance
(ElEmEr)} names the proposed underlying dynamical field/state structure.  ``Resonance''
does not mean that every \ElEmEr{} state is already a closed resonance.  An open travelling
mode, a recurrent closure, and a more extended organization may all be different states
of the same underlying field.

The picture is deliberately field-first.  What persists fundamentally is not a set of
immutable labelled parcels but the evolving local field state.  An open organization can
carry a disturbance across space; a recurrent one can furnish a physical cycle; a
localized recurrent one can be particle-like; and interacting organizations can alter
the complete local state from which subsequent propagation continues.  Propagation,
clocks, interaction, and closure are therefore different aspects of one dynamical
problem.  An \emph{ElEmEr wave} is a propagating or standing wave-like pattern of field
excitation, not a separately conserved piece of substance; an \emph{ElEmEr field excitation} more generally covers organized disturbances that need not be wave-like.  Individual field-variable values need not by
themselves specify the complete local dynamical state.

Time is not taken to be a fourth material or spatial-like dimension of this primitive
ontology.  Four-vectors, null cones, and Lorentz tensors remain available when an
effective relativistic description calls for them; using that language does not make
material four-dimensional spacetime part of the primitive ontology.

Spatial arclength, speed, orthogonality, gradients, and transverse projection require
local spatial geometry.  We write the local line element as
$\dd\ell^2=h_{ij}\,\dd x^i\dd x^j$, where $h$ is a positive-definite
spatial metric on a working three-space chart.  The three-vector formulas below use
an oriented local orthonormal frame (so $h_{ij}=\delta_{ij}$ in that frame), with
metric-compatible spatial differentiation where required.  These are explicit
kinematic/mathematical conventions, not a derivation of one unique global background
geometry or of a frameless ground state; that physical completion remains open.

\subsection{Primitive succession and local-\texorpdfstring{$c$}{c} propagation}
\label{subsec:succession}

Primitive change propagates locally at $c$ through physical three-space, while $s$
provides the primitive before--after ordering of successive physical states.  Along a
resolved primitive \ElEmEr{} propagation characteristic,
\begin{equation}
  \boxed{\dd\ell=c\,\dd s,}
  \label{eq:dlds}
\end{equation}
with $\dd\ell$ the physical spatial arclength.  The equation is local: advancing through
one increment of primitive succession advances the propagation characteristic through a
corresponding spatial arclength at speed $c$.

The role of $s$ stops there.  It is not an observer-readable master clock, and the local
relation above does not tell us how independently described histories acquire a common
origin.  A single globally aligned $s$ is therefore a candidate foundational possibility,
not a consequence already obtained from local-$c$ propagation.

Continuum field equations involving $\partial_s$ and neighbouring spatial points
require a compatible \emph{local evolution chart} in the domain where those equations
are used.  This local-chart compatibility is an additional working antecedent, not a
construction of globally aligned equal-$s$ slices, their physical synchronizability,
or a globally privileged frame.  A native law must ultimately state how local succession
increments and the complete field state fit together across interacting characteristics;
$\dd\ell=c\,\dd s$ alone does not resolve that compatibility.

By \emph{local-$c$ propagation} we mean that the resolved primitive propagation tangent
has magnitude $c$, independent of its propagation direction and source motion.  This
local tangent describes the field evolution itself; it is not the centre-of-mass velocity
of a closure or the bulk velocity of a material medium.  Its magnitude can remain $c$ while direction,
phase, internal state, mode content, and larger-scale organization change.

\subsection{Recurrent clocks and operational time}
\label{subsec:operational_time}

A physical clock is a repeatable process.  Recurrence therefore supplies the natural
bridge from primitive succession to measured time: a clock reading can be a cycle count
or calibrated phase of a recurrent organization, while observer time is constructed by
comparing clocks and exchanging physical signals.

The three notions are kept separate:
\begin{equation}
\boxed{
\begin{array}{rcl}
 s &:& \text{primitive succession/evolution order},\\
 \tau_C &:& \text{local recurrent-process or closure-clock reading},\\
 t_{\mathcal O} &:& \text{observer/network coordinate time built from clocks and signals}.
\end{array}}
\label{eq:three_times}
\end{equation}
A moving or externally perturbed clock may change its cycle geometry and therefore its
measured rate while every resolved primitive segment still propagates locally at $c$.
Relativistic clock behaviour could therefore arise from a changed recurrent process rather
than a changed primitive propagation speed; the Lorentz-form completion is discussed in
\cref{sec:lorentz_effective}.

\subsection{Open propagation and recurrent organization}
\label{subsec:open_closed_premise}

The same field can continue outward as an open mode or bend back into a compatible
return.  If redirection, phase, and internal state repeatedly reproduce a compatible
configuration, propagation becomes recurrent.  A \emph{closure} is a sufficiently
localized and dynamically sustained recurrent organization of the \ElEmEr{} field.

The constructive picture is
\begin{equation}
\boxed{
\begin{gathered}
\text{open propagation}
\longrightarrow
\text{deflection / return}
\longrightarrow
\text{recurrent organization}\\
\longrightarrow
\text{persistent localized state}.
\end{gathered}}
\label{eq:closure_story}
\end{equation}
Localization need not be produced by lowering the primitive propagation speed.  The
field may remain locally luminal while repeated turning and return build a finite
organization.  Geometric return alone, however, cannot tell us which configuration is
actually captured or sustained; that requires the dynamics of the field itself.

\section{From propagation to local response}
\label{sec:structural}

A travelling disturbance need not keep the identity of a little object while
passing through an interaction.  What matters is how the complete local field
configuration changes and what it sends onward.  This gives the open-mode,
closure, and clock picture a common physical interaction language: the state
present at a location determines which responses are possible, while a complete
native law would determine which actually occur.

\subsection{The complete local field state is the interaction object}
\label{subsec:field_first}

At a physical location there is one current \ElEmEr{} state.  Labels such as ``incoming
mode A'' and ``incoming mode B'' may describe how a configuration was prepared, but they
need not remain primitive identities after overlap, redirection, or mode conversion.
The local dynamics acts on the state that is actually there.  This is an
account of the interaction object, not yet a rule for measurement outcomes
in entangled preparations (see \cref{sec:companions}).

Accordingly, the interaction problem is not fundamentally ``what force does particle A
exert on particle B?'' but ``how does the combined local field state reorganize its
momentum, stress, phase, and internal mode content?''  Schematically,
\begin{equation}
\boxed{
\begin{gathered}
\text{complete local state}
\longrightarrow
\text{local reorganization}\\
\longrightarrow
\text{changed propagation / steering / mode structure}.
\end{gathered}}
\label{eq:field_first_chain}
\end{equation}
Particle language can still be used for stable incoming and outgoing organizations; it
is not required as the microscopic ontology of the overlap itself.

\subsection{A candidate interaction--response architecture}
\label{subsec:interaction_response}

We use \emph{interaction opportunity} for a local state that has departed from the
locally valid free-propagation class and therefore admits the possibility of an
interaction-induced response.  Departure from the free class is not by itself a
sufficient activation law.  The response, if a physical
channel activates, is a reorganization of local field, stress, and current content:
\begin{equation}
\boxed{
\begin{gathered}
\text{sector source / interaction condition}
\longrightarrow
\text{local field--stress--current reorganization}\\
\longrightarrow
\text{steering / state / mode response at local }c.
\end{gathered}}
\label{eq:response_architecture}
\end{equation}

The same architecture can accommodate physically different outcomes.  Momentum
redistribution may bend a propagating mode; internal structure may rotate or change
mode; a sufficiently organized response may help form a recurrent return; repeated
interactions may redistribute spectral or occupation content.  These outcomes need not
share one microscopic force law.  The common element is the local reorganization of the
complete state.

Leaving the free class is only an opportunity for interaction.  Non-collinearity,
counterflow, or other representation-level witnesses can diagnose such a departure, but
the activation condition, sign, strength, range, and constitutive law belong to the
specific interaction dynamics.

\subsection{Generation and causal field extension}
\label{subsec:causal_propagation}

The candidate causal picture separates production of a change from transport of that
change.  A local interaction may first alter the physical field/state.  Once generated,
that alteration propagates through neighbouring field states along successor-directed
local-$c$ characteristics:
\begin{quote}
A local interaction may generate a physical field/state change; the generated change
then propagates successor-directed along primitive local-$c$ characteristics.
\end{quote}

We refer to this as \emph{causal field extension / state propagation}: a disturbance
produced here reaches later locations by propagation rather than by an instantaneous
update of a spatially extended stress network.  ``Successor-directed'' refers to the local causal
before--after relation along the characteristic; it does not assume a universal spatial
centre or a derived global synchronization of $s$.

\subsection{A candidate causal-reachability constraint}
\label{subsec:causal_reachability}

A change that has propagated away from a source cannot simply be recalled by
reversing that source later.  The later influence must itself travel through
space at the local propagation speed.  To recreate the full reversed process
would require suitably prepared incoming field data, including directions,
phases, internal state, and correlations; it is not equivalent to turning the
visible source around.  In that sense a present configuration may carry
consequences of interactions that are no longer local to the original source.

This provides a candidate distinction between reversing a local preparation
and reversing a complete physical history.  It is not a native proof of
irreversibility, an entropy law, or a prohibition on mathematical reversed
solutions.  The precise complete-state reversal test and its reachability
limitations are given in \cref{subsec:complete-state-reversal}.


\subsection{Steering while primitive speed remains fixed}
\label{subsec:projected_steering}

We use \emph{steering} for a local reorientation of resolved propagation while its
primitive speed remains fixed. When a ray-like field component bends, its propagation
direction changes while the magnitude of its primitive motion remains fixed at $c$.
A succession of such turns is also the geometric ingredient a recurrent organization
would need. The bulk motion of a whole recurrent pattern is a separate quantity: it can
change even though its resolved primitive tangents remain locally luminal.
A fixed propagation magnitude therefore makes transverse response, rather
than microscopic slowing-down, the natural steering geometry.  In a three-vector response description, let $\mathbf u$ denote such a
resolved local propagation tangent with $|\mathbf u|=c$.  A speed-preserving steering law
can be written
\begin{equation}
\boxed{
\frac{\dd\mathbf u}{\dd s}
=
\mathbf P_\perp(\mathbf u)\,\mathbf g_{\rm st}
+
\frac{\mathbf u}{c}\times\mathbf h_{\rm st},
}
\label{eq:projected_steering}
\end{equation}
with
\begin{equation}
\mathbf P_\perp(\mathbf u)=
\mathbf I-\frac{\mathbf u\otimes\mathbf u}{c^2}.
\end{equation}
Both terms are transverse to $\mathbf u$, so the equation bends or reorients the
propagation without changing its primitive speed.  The speed-preservation check is given in \cref{subsec:speed-proof}.
The geometry is fixed at this level;
the microscopic source and the physical meaning of $\mathbf g_{\rm st}$ and
$\mathbf h_{\rm st}$ depend on the interaction law of the sector being described.

An effective Lorentzian representation yields a related conditional algebraic result.
Let $g$ be the effective Lorentzian metric and $A$ one common real \emph{infinitesimal}
linear operator for all null directions at a point.  Require tangency to the null cone:
$g(K,AK)=0$ for every nonzero $K$ satisfying $g(K,K)=0$.  In Lorentzian
dimension at least two, this implies
\begin{equation}
\boxed{A=F+\phi I,\qquad F_{\mu\nu}=-F_{\nu\mu},}
\label{eq:all_null_A}
\end{equation}
where $F_{\mu\nu}=g_{\mu\alpha}F^\alpha{}_{\nu}$.  The scalar part generates
uniform rescaling of the null vector.  The Lorentz generator $F$ can change its
projective propagation direction and also rescale its normalization in a
direction-dependent way; in $1+1$ dimensions its action preserves each null
line.  A null-vector normalization is not, without a specified physical
calibration, a field amplitude or an independently established frequency.
The proof, including the $1+1$ case, is given in \cref{subsec:all-null-proof}.  This is an \emph{infinitesimal} cone-tangency condition;
it does not assert that $AK$ itself is null.  Two-form-like steering is natural within
this specified local-linear representation, but the theorem does not derive a unique
microscopic physical two-form or its source law.
For a specified, physically normalized, direction-dependent transverse response,
the all-directions test of this \emph{single common operator} is summarized
in \cref{subsec:sphere_response_test}.  Agreement for one null direction
alone is insufficient.

\section{Closure as persistent physical organization}
\label{sec:closure_gates}

Closure offers a way to understand persistence without postulating a second
primitive substance.  An open mode carries change away; a returning mode revisits
its own organized region; a captured, sustained return could keep a localized
identity over many cycles.  In that picture a particle-like object is not a
bead moving around a wire.  What persists is the field's continuing organized
motion and the relations that make each return compatible with the next.
The present framework has consistency requirements for return, but not yet
the complete dynamical capture mechanism.

\subsection{Open, closed, and extended organizations}
\label{subsec:taxonomy}

Three broad organizational descriptions are useful:
\begin{description}[style=nextline,leftmargin=2.3cm]
  \item[Open] a propagating component that does not itself form an autonomous recurrent return;
  \item[Closed] a localized recurrent organization with a proper-return cycle;
  \item[Extended] a causally prepared organization whose relevant state spans separated
  regions without requiring instantaneous signalling.
\end{description}
In the intended physical reading, photon-like radiation is associated with open
propagating modes, not with microscopic permanent beads of substance.  A stable
matter-like object would instead require a dynamically sustained recurrent
organization.  An extended preparation could relate several regions through
its physical state and causal history without becoming a new primitive
substance.  These are ways of describing organization, not a completed native
particle identification or measurement theory.

These are neither a demonstrated exhaustive partition of all field states nor mutually
exclusive labels for a complete physical configuration.  An extended organization may
contain localized closures, while the same complete field state also includes open
propagating modes.  The labels describe organizational features, not primitive substances
or particle species; spin, charge, family, and detailed stability require further structure.

For a closure, the path is not a pre-existing wire and the field is not a bead travelling
around it.  Path, phase, stress, and internal state are properties of the recurrent
configuration.  In a \emph{simple scalar spatial-phase return model}, a consistency
condition can be written
\begin{equation}
\oint_\Gamma k_{\rm eff}\,\dd\ell=2\pi n,
\label{eq:return_condition}
\end{equation}
where $k_{\rm eff}$ denotes a model-specific \emph{effective return-phase rate}, not an
assertion that a primitive travelling crest accumulates phase along its own null
characteristic.  The full return phase can also contain geometric, interaction, and
internal-frame contributions.  For a specified multicomponent transport model, an
amplitude-level return criterion has the more general schematic form
\begin{equation}
 e^{\mathrm i\Theta_\Gamma}\,\mathcal H_\Gamma\Psi_0=\Psi_0,
 \label{eq:general_return_holonomy}
\end{equation}
where $\Theta_\Gamma$ collects the scalar phase contributions not already contained in
the internal transport/holonomy operator $\mathcal H_\Gamma$.  Whether exact amplitude
return or only physical-ray return is appropriate must be stated for each model;
nontrivial frame holonomy is not selected merely by drawing a closed path.  These are
\emph{consistency conditions}, not capture or binding mechanisms.  A physical closure
must also acquire amplitude support, enter the recurrent state dynamically, survive
perturbations, and control whatever leakage remains.

The older C-motion picture emphasized that a closure is \emph{sustained by organized
motion}.  That intuition survives the present status discipline in a narrower form:
recirculation, internal steering, mode conversion, and finite-$c$ self-consistency are
natural ingredients to investigate, but none is presently established as the universal
microscopic binder.

\subsection{From interaction opportunity to persistent closure}
\label{subsec:return_not_binding}

Formation separates into several physically different stages.  A local state must make
a response channel available; the response must steer or convert the field into a
return-compatible geometry; phase and spectral relations must support recurrence; excess
energy or momentum may need somewhere to go; and the resulting configuration must remain
stable and sufficiently isolated.

A successful formation sequence may therefore involve:
\begin{enumerate}[label=(\roman*)]
  \item interaction opportunity;
  \item activation of a physical response channel;
  \item steering and geometric/polarization compatibility;
  \item phase/spectral capture;
  \item mode conversion where required;
  \item nonlinear capture into a recurrent organization;
  \item energy/momentum shedding or relaxation where required;
  \item finite-radius dynamical stability;
  \item long-term isolation or sufficiently suppressed leakage.
\end{enumerate}

Calling all of these steps ``binding'' hides the main physical questions.  Current results
constrain parts of the sequence, especially the distinction between steering, return, and
capture; the microscopic law that carries a configuration through the full sequence is
still unknown.

\subsection{Rest-like mass from proper-return energy}
\label{subsec:mass_inertia}

A finite recurrent organization may carry invariant energy even though every
resolved primitive propagation segment remains locally at $c$.  To make the
mass interpretation precise, first assume an \emph{isolated, self-consistent}
recurrent field organization for which the effective dynamics supplies a
conserved total energy $E$, a conserved total momentum $\mathbf P$, and a
physical rest frame.  We use \emph{proper-return energy} to mean the energy
$E_{\rm return}:=E_{\rm rest}$ of the \emph{complete} organization in that
rest frame.  ``Return'' describes the recurrent structure; $E_{\rm return}$
is not energy integrated over a period, the energy of one carrier strand, or a
frequency-to-energy formula.

Within the further \emph{Conditional} effective Lorentz-covariant
energy--momentum description,
\begin{equation}
  E^2-c^2|\mathbf P|^2=M^2c^4,
  \qquad
  \boxed{Mc^2=E_{\rm return}\quad(\mathbf P=0).}
  \label{eq:mass_return}
\end{equation}
This rest-frame identity defines/identifies the effective invariant-mass
parameter once the total energy--momentum of a recurrent solution is known;
it does not derive that solution, its rest energy, or its absolute mass from
primitive \ElEmEr{} dynamics.  The mass-like property belongs to the
organization as a whole, not to a primitive propagating component that has
slowed below $c$.

The same picture motivates a \emph{candidate} account of inertia.  Changing
the collective motion of a closure may require changing the geometry and
energy--momentum flow of its entire return pattern; a rigid translation of an
unchanged internal loop is not established as a valid native solution.  If
the effective Lorentz-covariant relation above holds for a positive-mass
branch, its small-momentum expansion is
$E(\mathbf P)=Mc^2+|\mathbf P|^2/(2M)+O(|\mathbf P|^4)$, supplying the
usual inertial coefficient \emph{within that effective description}.
This is not yet the microscopic response of a boosted native closure.
A native derivation would require the completed action, a self-consistent
moving localized solution, and its computed conserved energy--momentum.
The conditional equality of rest and inertial mass must not be presented as
a derivation of universal gravitational free fall or the equivalence principle.

The need to account for the \emph{whole} organization is not merely verbal.
Under specified isolation and conservation assumptions, the Laue/virial
identity \cite{Laue1911Dynamik,Giulini2018LaueRevisited} ties a positive rest energy to a nonzero integrated trace of the
complete stress tensor.  A purely traceless complete, isolated stationary
stress system cannot furnish positive rest energy in that class.  The
assumptions, integral identity, and exceptions are given in
\cref{subsec:stress-trace-mass}.  This is a constraint on any proposed binder,
not a calculation of an electron or any other closure mass.


\subsection{What must sustain a closure?}
\label{subsec:binder}

Geometric return says what a recurrent state must look like after a cycle; it does not say
what captures the field into that cycle or keeps the radius finite.  We use
\emph{binder} as shorthand for whatever microscopic interaction or self-organization
mechanism supplies that capture and sustained finite-radius localization.  Several simple
possibilities are already constrained.  With the measured weak-field Newtonian
coupling, ordinary gravity is much too weak to supply elementary-particle-scale
self-localization: the compactness $GE/(Rc^4)$ is tiny for the particle energies
and radii at issue.  This is not an exclusion of strong gravity at arbitrary
energies or radii.  The standard positive-energy, same-charge Abelian vector
branch likewise has repulsive classical self-energy and does not provide a
direct contractile binder without additional stresses or interactions.
The historical electron self-stress problem already required additional
cohesive stress \cite{Poincare1906Electron}; that comparison is not a theorem
excluding arbitrary nonlinear or additional-field models.  This restricted negative control
does not exclude a distinct fermion--antifermion or left/right coherence-channel
instability.  Conversely, an attractive susceptibility or coherence instability
would not by itself establish physical capture, full stress balance, or stability.

One possibility considered in the framework is an additional strong-like microscopic
interaction within the \ElEmEr{} dynamics.  Its appeal is physical rather than formal: a
short-range nonlinear channel could, in principle, supply capture and finite-radius
support without being identified with the weak long-range steering already used in other
sectors.  No particular hidden gauge group, coupling, or action has been established.
Any successful mechanism must produce capture, finite radius, consistent energy--momentum
exchange, and a stable or sufficiently long-lived recurrent state.

\section{Energy, action, stress, and dynamical response}
\label{sec:energy_action}

A field pattern can move locally at the same primitive speed as another
pattern yet carry different energy, momentum, internal organization, and
stresses.  Thus the fixed-speed premise is not an energy formula.  What
holds a recurrent state together, what changes its collective motion, and
what its clock measures must be connected by the actual dynamics.
Energetics enter when that dynamics supplies an action or Hamiltonian
structure.  From that structure come energy and momentum for organized
states, stresses and fluxes for their local transport, and ultimately the forces or
steering responses generated by redistribution of that dynamical content:
\begin{equation}
\boxed{
\begin{gathered}
\text{action / evolution structure}
\longrightarrow
\text{energy, momentum, stress, flux}\\
\longrightarrow
\text{effective response and steering}.
\end{gathered}}
\label{eq:energy_hierarchy}
\end{equation}
The same local-$c$ kinematics can therefore support physically different energy-bearing
organizations depending on the action and the state.

\subsection{Energy and action charge are different quantities}
\label{subsec:energy_action_layers}

Where an autonomous action or Hamiltonian exists, energy is its evolution/translation
charge in the usual Noether/Hamiltonian sense \cite{Noether1918,PeskinSchroeder1995QFT}.
A different conserved quantity can arise when the action has an exact continuous phase
symmetry.  We call the corresponding Noether charge $Q$ the \emph{action charge}.
Here ``charge'' means a conserved quantity generated by a symmetry; it is not electric
charge, gauge charge, or particle number.  In the normalization used here, $J_0Q$
carries the corresponding action scale, with $J_0$ supplying the overall action
normalization.  Energy and action charge can therefore be conjugate within one dynamics
without being the same conserved quantity.  There is also an irreducible
\emph{classical normalization freedom}: replacing a complete classical action
$S$ by $\lambda S$ with $\lambda>0$ leaves its Euler--Lagrange equations and
trajectories unchanged, whereas action-derived currents, energies, stresses and
symplectic form scale with $\lambda$.  A single shared action may relate
relative coefficients across sectors but cannot, from those classical equations
alone, fix its remaining overall scale.  Consequently the physical assignment
$J_0=\hbar$ needs an independent normalization/quantization or empirical input;
it is not selected by writing down a common action.

Assume an autonomous differentiable action/Hamiltonian with finite energy and
an exact continuous phase charge, admissible boundary conditions, and a smooth
family of stationary relative equilibria (or recurrent states for which the
same constrained-stationarity condition is established).  At each family member,
require $\delta(E-J_0\Omega Q)=0$ with $\Omega$ held fixed in this field variation
and all remaining charges and topological data fixed along the family.
Evaluating that variation on the family tangent gives
\begin{equation}
\boxed{\frac{\dd E}{\dd Q}=J_0\Omega.}
\label{eq:conjugacy}
\end{equation}
Here $J_0$ is the overall action-normalization scale, so that $J_0Q$ has action
dimensions when the symmetry charge $Q$ is dimensionless.  The meaning of $\Omega$
depends on the evolution convention of the assumed autonomous effective action: if
$\sigma$ is the parameter whose translations are generated by $E$, and $\vartheta$
is the recurrent symmetry phase, then $\Omega=\dd\vartheta/\dd\sigma$ in that same
convention.  A compatible local succession chart might use $\sigma=s$, but this
requires its own Hamiltonian and normalization antecedents.  A frequency measured
against a closure clock $\tau_C$ or observer time $t_{\mathcal O}$ cannot be substituted
without an additional dynamical/operational conversion; nor is the native physical
energy generator already known.  The relation is a Conditional result within a
specified recurrent symmetry family, not a universal energy formula for arbitrary
ElEmEr waves or closures.  It leaves the numerical value of $J_0$ undetermined;
the empirical identification $J_0=\hbar$ and the associated $E=\hbar\omega=hf$
bridge are treated in Quantum Foundations.
This is the constrained stationary-family relation familiar from Q-ball/soliton
analysis \cite{Coleman1985Qballs}, with the explicit normalization $J_0$ retained
here.  It neither constructs a native C-motion action nor selects a stable
closure.  The short family-tangent argument is given in
\cref{subsec:conjugacy-proof}.

For a nonlinear family the differential relation does not generally integrate to
\begin{equation}
E=J_0\Omega Q.
\end{equation}
Instead,
\begin{equation}
\boxed{E(Q)-E(Q_0)=J_0\int_{Q_0}^{Q}\Omega(q)\,\dd q,}
\label{eq:energy_integral}
\end{equation}
where the reference state fixes the integration constant.
The simple product identity $E=J_0\Omega(Q)Q$ on a smooth interval with $Q\ne0$
would additionally require $Q\,\dd\Omega/\dd Q=0$ by differentiation and therefore
constant $\Omega$ on that interval, with the energy zero selected accordingly.
A product relation found at one particular state does not establish an identity for
the full nonlinear family.

A related \emph{Open (sharp)} question is whether a generic finite wave admits a
physically normalized, positive, conserved local action current with a justified
mode decomposition.  A charge defined in a specified phase-symmetric effective
action, or a factorization for a resolved mode, does not establish this for an
arbitrary coherent/interacting field.  The missing native current and its
normalization remain separate from the stationary-family identity; the quantum
normalization bridge belongs to Quantum Foundations.

\subsection{Stress redistribution and steering}
\label{subsec:stress_steering}

A field-first interaction must also account for energy and momentum locally.  For
\emph{one resolved, locally one-way null-flow branch} with $\varepsilon>0$ and
$|\mathbf u|=c$, assume the null-flow calibration
$\mathbf p_{\rm car}=\varepsilon\mathbf u/c^2$ and
$\boldsymbol\Pi_{\rm car}=\varepsilon\,\mathbf u\otimes\mathbf u/c^2$.
In a compatible local evolution chart, with the spatial metric/volume element held
fixed for this balance or their evolution terms included in the exchange terms,
write the carrier--interaction exchange equations as
\begin{align}
\partial_s\varepsilon+\nabla\!\cdot(\varepsilon\mathbf u)&=q,
\label{eq:resolved_null_energy_balance}\\
\partial_s\!\left(\frac{\varepsilon\mathbf u}{c^2}\right)
+\nabla\!\cdot\!\left(\frac{\varepsilon\mathbf u\otimes\mathbf u}{c^2}\right)
&=\mathbf f,
\label{eq:resolved_null_momentum_balance}
\end{align}
where $q$ and $\mathbf f$ are the exchanged power and force densities, and
$D_s:=\partial_s+\mathbf u\!\cdot\nabla$ denotes transport along the resolved branch.
These assumptions give the transverse relation
\begin{equation}
\boxed{
 D_s\mathbf u=\frac{c^2}{\varepsilon}\,
 \mathbf P_\perp(\mathbf u)\mathbf f.}
\label{eq:force_balance}
\end{equation}
The longitudinal part yields the compatibility condition $q=\mathbf u\cdot\mathbf f$
(\cref{subsec:null-flow-proof}).  This describes energy exchange in the assumed
calibrated null-flow model; it is not an identification of that exchange with
null-vector rescaling in \cref{eq:all_null_A}, nor with physical amplitude change.
For a locally conserved total momentum with an interaction-sector momentum density
$\mathbf p_{\rm int}$ and stress flux $\boldsymbol\Sigma_{\rm int}$,
$\mathbf f=-\partial_s\mathbf p_{\rm int}-\nabla\!\cdot\boldsymbol\Sigma_{\rm int}$.
Only when the interaction momentum-density derivative is negligible, stationary, or
already included in an effective flux may this reduce to the stress-only expression
\begin{equation}
\boxed{
 D_s\mathbf u=-\frac{c^2}{\varepsilon}\,
 \mathbf P_\perp(\mathbf u)\!\bigl(\nabla\!\cdot\!\boldsymbol\Sigma_{\rm int}\bigr).}
\label{eq:stress_divergence_steering}
\end{equation}
Thus a stress imbalance can steer a \emph{resolved} null flow without changing its
primitive speed.  A generic coherent superposition, counter-propagating pair, or full
interacting field need not have the rank-one null-flow stress above.  Neither the
carrier--interaction split nor the source/stress law is a derived universal ElEmEr
constitutive equation; those require native dynamics.

\section{Operational relativity from physical clocks and signals}
\label{sec:lorentz_effective}

Relativity enters \cmotion{} through physical measuring systems.  Rods, recurrent clocks,
and signalling procedures are themselves field organizations built from the same
local-$c$ propagation.  Their mutual relations can therefore acquire Lorentz form without
requiring primitive time to be a material coordinate on the same footing as the three
spatial directions.

The recurrent-clock picture gives a concrete candidate mechanism.  A moving clock is a
returning field pattern with collective drift.  If that drift changes the spatial and
phase geometry through which the pattern must complete its return, the cycle rate can
change even though each resolved primitive segment still propagates locally at $c$.
Older simple perpendicular-motion constructions capture this intuition but are not a
general derivation; an orientation-independent moving-closure calculation is still
needed.

At the effective level, when the standard Lorentz relations apply
\cite{Einstein1905Relativity},
\begin{equation}
\dd\tau=\frac{\dd t}{\gamma},\qquad
\gamma=(1-v^2/c^2)^{-1/2}.
\label{eq:effective_lorentz_clock}
\end{equation}
Here $\tau$ is a physical clock reading and $t$ an observer coordinate, not a direct
measurement of primitive succession $s$.

A complete microscopic account must recover the same operational relations for clocks,
rods, signalling, and interactions.  If a single global alignment of $s$ exists, that
structure must not produce an observable absolute-motion dependence.  This no-$s$-leak
requirement and the general moving-clock/rod construction remain open; local $c$ alone is
not being used as a proof of Lorentz covariance.  Pure open-ray
transparency, in particular, does not ensure boost-compatible behaviour
for interacting mixed states; a restricted countermodel is recorded in
\cref{subsec:no-s-leak-countercontrol}.
\section{Scale, discreteness, and structural constraints}
\label{sec:constraints}

The geometry of recurrence places real constraints on a completion even when it does not
supply the missing microscopic law.  In particular, it separates questions of mode
selection from questions of absolute scale and rules out several tempting shortcuts.

\subsection{Absolute scale requires an absolute source}
\label{subsec:scale_source}

A scale-covariant closure problem may determine shapes, ratios, winding conditions, or
mode labels without determining whether the physical object is a femtometre or a metre
across.  An absolute ruler requires a dimensionful input, a boundary or state scale, or a
dynamical mechanism such as anomaly-driven transmutation.  Established examples
include radiative scale generation in the Coleman--Weinberg construction \cite{ColemanWeinberg1973}
and running asymptotically free non-Abelian theories
\cite{GrossWilczek1973,Politzer1973}.  These are examples in specified quantum
field theories, not a demonstrated scale-generation mechanism for native
\ElEmEr{} dynamics.

Topology, return phase, non-radiation, or retardation may therefore help decide
\emph{which} recurrent mode is allowed while leaving its overall size undetermined.  Until
a dynamical scale is obtained, the present work-in-progress formulation allows at most one common energy scale
\begin{equation}
\boxed{\Lambda_C}
\end{equation}
as a single framework-wide fallback absolute scale.  It is not a separate fitted scale
for each particle or sector and is used only until a dynamical scale-generation mechanism
is established.

The overall action normalization $J_0$ is a separate physical question.  An energy scale
sets an absolute energetic ruler; an action normalization sets the relation between
phase/action charge and energy.  They should not be hidden inside unrelated fit
parameters.  If $J_0$ has action dimensions and the energy/frequency conventions are
consistently fixed, one may define dimensional \emph{reference} scales
\begin{equation}
 \ell_C:=\frac{J_0c}{\Lambda_C},\qquad
 s_C:=\frac{J_0}{\Lambda_C}.
 \label{eq:scale_reference_dimensions}
\end{equation}
Here $s_C$ is a reference \emph{succession interval}, not a directly readable clock
period.  Neither quantity is a predicted closure radius, internal wavelength, or
recurrence period: those require a physical closure law and dimensionless state data.
Conversely, $c$ and an energy-valued $\Lambda_C$ alone cannot furnish an absolute length
without an independently justified action normalization or equivalent dimensional
input.  Identifying $J_0$ empirically with $\hbar$ does not derive that normalization
from ElEmEr dynamics and remains the separate Quantum Foundations bridge.  The
one-scale fallback permits one common energy calibration, not independent fits of
$J_0$, closure radii, or sector-specific microscopic lengths.

\subsection{Return constraints can discretize without explaining the whole spectrum}
\label{subsec:return_discreteness}

A recurrent boundary-value problem can admit only certain phase-compatible returns.  In
that limited but important sense, continuous local propagation can organize into a
discrete set of allowed recurrent modes.  The return condition supplies spectral
selection only after the closure dynamics has been specified; identifying those modes
with the observed particle spectrum additionally requires the microscopic action, scale,
and stability structure.

\subsection{No common nonzero additive count per outgoing photon}
\label{subsec:additive_no_go}

Photon multiplicity in positronium annihilation provides a scoped test of an
identical, exactly conserved, nonzero additive weight per outgoing $\gamma$
quantum.  Charge-conjugation conservation permits odd photon multiplicities for
ortho-positronium and even multiplicities for para-positronium; permission alone
does not establish a nonzero rate.  QED calculations predict rare five-photon
and four-photon channels in addition to the dominant three-photon and two-photon
channels \cite{AdkinsBrown1983FivePhoton,Lepage1983Multiphoton}.
Reported five-photon branching measurements are consistent with QED, but their
statistical precision does not independently establish a nonzero branching ratio
against zero at conventional discovery significance
\cite{Matsumoto1996FivePhoton,Vetter2002FivePhoton}.
They must not be described as a decisive experimental observation of nonzero
five-photon decay.  The para-positronium $2\gamma/4\gamma$ comparison supplies
stronger direct experimental support for differing multiplicities from the same
initial-state class \cite{Vetter2002FivePhoton}.

Conditional on the nonzero $3\gamma$ and $5\gamma$ channels predicted by QED,
let the same prepared ortho-positronium state carry a fixed quantity
$W_{o\mathrm{Ps}}$.  If that quantity is exactly additive and conserved, every
outgoing photon carries the same state-independent constant weight $w_\gamma$,
and no other outgoing carrier or environment carries it, then
\begin{equation}
\boxed{W_{o\mathrm{Ps}}=3w_\gamma=5w_\gamma
\quad\Longrightarrow\quad w_\gamma=0.}
\label{eq:fixed_photon_count_no_go}
\end{equation}
The experimentally better-supported para-positronium comparison gives the same
algebraic obstruction with $2w_\gamma=4w_\gamma$, under the corresponding fixed
initial-state and complete-accounting assumptions.  Thus the Negative conclusion
excludes a universal fixed nonzero photon-count action assignment under these
explicit premises, not every possible conserved action structure.  It does not
exclude energy--momentum, signed charges, state-dependent weights, continuous
field-level measures, or nonadditive global quantities.

Here the observations count outgoing $\gamma$ quanta, not a representation-dependent
number of field modes.  Photon language describes measured final states; it does
not make photons primitive material corpuscles in \cmotion{}.  A genuine action
or closure quantum requires a specified conserved or normalized structure, not
an imposed universal tally of permanent parcels.

\paragraph{Field-first interpretation of the two controls.}
The positronium multiplicities give the sharper same-initial-state fixed-count
obstruction; a separate charged-lepton conversion/decay illustration reinforces why
energy-dependent changes in outgoing field organizations cannot automatically be understood
by tracking permanently labelled wave strands (details in
\cref{subsec:conversion-cycle}).  This motivates the field-first
ontology adopted here: propagating waves are real \emph{organizations or modes}
of the ElEmEr field, while their mode number and apparent strand identity need
not be conserved.  It does \emph{not} assert that wave dynamics is impossible,
that waves have no physical reality, or that experiment has uniquely proved a
particular microscopic ElEmEr field law.

\subsection{The observed three-family count remains unexplained}
\label{subsec:family_concession}

\cmotion{} does not presently derive why there are exactly three fermion families.
Radial excitation labels, scale or coupling ladders, and simple branch counts have not
forced the observed number.  If family structure is ultimately native to the framework,
it may require a deeper internal, topological, or nonlinear organization of
closures.  In particular, irreducible Clifford/Weyl module dimensions do not
supply the observed count of three; the bounded exclusion and its assumptions
are discussed in the Quantum Foundations companion
\cite{CmotionQuantumFoundations2026}.  Other tested family routes remain
bounded exclusions or open research, not a native family-count derivation.

\subsection{Global phase and local gauge transport are different layers}
\label{subsec:phase_gauge}

An exact global phase symmetry can supply a Noether charge once the corresponding action
is specified.  Turning such a symmetry into local frame or gauge transport is a further
physical construction requiring a local frame structure, connection, and sector
dynamics.  Global phase structure can therefore play a role in the energy/action layer
without already supplying local $\mathrm U(1)$, $\mathrm{SU}(2)$, Maxwell, or other gauge
dynamics.

\section{Where the construction continues}
\label{sec:companions}

The present paper stops at the common architecture.  The companion manuscripts
address additional assumptions and calculations without folding those results
back into the foundational premises.  The planned first wave comprises this
Main Framework and five companions: Quantum Foundations, Thermodynamics, Gravity,
Structural Entanglement I, and Gauge Architecture and Constraints I.  These are
ongoing manuscripts, not assertions of already published results.  Electroweak
Modes, Standard Model Structure II, and the separate measurement-completion
programme remain deferred research:

\begin{longtable}{@{}P{4.0cm}P{10.6cm}@{}}
\toprule
Topic & Developed separately \\
\midrule
\endfirsthead
\toprule
Topic & Developed separately \\
\midrule
\endhead
Quantum foundations \cite{CmotionQuantumFoundations2026}
& The status of $J_0=\hbar$ and $E=\hbar\omega=hf$, Clifford/spinorial structure, and other quantum-foundation questions. \\
Thermodynamics \cite{CmotionThermodynamics2026}
& Entropy, temperature, equilibration, thermalization, and the thermodynamic arrow. \\
Gravity \cite{CmotionWeakFieldGravity2026}
& The gravitational source and coupling, conditional rank-2 weak-field route, PPN recovery, and strong-field questions. \\
Gauge Architecture and Constraints in \cmotion{} I \cite{CmotionSMStructureI2026}
& Internal-frame and gauge constraints and chirality;
\emph{C-Motion Standard Model Structure II: Confinement and Matter-Spectrum
Constraints} remains separately deferred research. \\
Electroweak modes (deferred research)
& B/W mode structure, SCF/Gap-C calculations, weak-angle bridges, and their numerical uncertainties; not a premise of this paper. \\
Entanglement \cite{CmotionEntanglementSinglet2026}
& Structural singlet correlations in the cited companion; measurement completion, Born-weight, and definite-outcome problems remain separately deferred research. \\
\bottomrule
\end{longtable}

\paragraph{Local causality is not yet a theory of quantum outcomes.}
Local causal propagation and a common-source extended organization are not,
by themselves, sufficient to establish the observed joint quantum outcome statistics.
A Bell-local factorized model with setting-independent preparation remains subject
to the CHSH bound even if its shared state is structurally elaborate
\cite{Bell1964,CHSH1969}.  Operational
no-signalling, Bell-local factorization, and spacelike operator microcausality are
distinct requirements.  Structural Entanglement I addresses singlet correlations;
the separately deferred measurement programme retains the Born-weight and
definite-outcome completion problems.

\subsection{A shared upstream dynamical problem}
\label{subsec:shared-native-dynamics}

Several questions addressed by these companion papers may depend on the same
microscopic information rather than requiring unrelated fundamental laws.
The current local-$c$ kinematics and field-first organization do not yet specify
a physically selected autonomous \ElEmEr{} action, or an equivalent complete
local evolution and constitutive law.  Such a law would have to determine how
the field responds to itself, how sources and interactions are activated, and
how energy and momentum are transferred between configurations.  Its existence
and physical selection remain open.

If that law were established, its equations and solutions could put several
otherwise separate problems on a common calculational footing.  Recurrent and
perturbed solutions would test binding, stability, characteristic scales and
internal spectra.  Its currents, stresses and interaction vertices could
constrain the effective quantum operator and action-charge structures,
gravitational sourcing and response, gauge and electroweak couplings, and the
physical preparation and analyzer interactions relevant to entanglement.
The accessible excitations and exchange processes would also be inputs to a
microscopic account of thermalization.  In this precise sense, the missing
native dynamics is a \emph{shared upstream bottleneck} rather than merely one
more unsolved item on a list.

This is a dependency map, \emph{not} an assertion that one guessed action will
solve all those problems.  An admissible action is not necessarily physically
selected, its solutions need not contain stable closures, and its symmetries
do not by themselves fix all constitutive coefficients.  Even a selected law
would still require separate demonstrations of action quantization, the Born
rule, thermodynamic equilibrium and a second law, gravitational equivalence,
and the observed particle spectrum and parameter values.  Those specific
completions remain with the relevant companion papers.  Constructing and selecting that native law lies outside the present paper;
it remains an open research problem.

The missing law would need to do more than recognize a possible interaction.
Given a complete local state, it would have to decide whether a source actually
activates, how the responding state changes, how strongly and quickly the change
proceeds, and---if more than one continuation is physically permitted---how the
continuation is selected or weighted.  Only then could one follow a proposed
history from an incoming field configuration through capture, free propagation,
or mode conversion and compare the result with an actual physical system.

The reduced native-dynamics selector tests explain why this task cannot be replaced
by checking the endpoints alone: different candidate laws can agree on the tested
free and recurrent endpoint properties while disagreeing on the intermediate
formation history.  These are scoped countermodels, not a theorem that nature
cannot select a unique law.  The precise diagnostic and conditional positive
selector control are collected in \cref{subsec:selector-controls}.


\section{Completion requirements and potential falsifiers}
\label{sec:falsifiability}

The story running through the paper is one proposed physical sequence:
propagating field change, locally activated reorganization, possible return
and sustained organization, and operational measurements made by those
organizations.  The aim is to explain why these aspects of physics might be
connected, rather than to identify them by analogy alone.  For that reason
the proposed physical picture is constrained by established observations and
must eventually make nontrivial, independently checkable predictions.  At
present, however, its native microscopic action and sector-specific laws are
not uniquely fixed.  The requirements below are therefore \emph{necessary
consistency and completion tests}, not five parameter-free predictions of the
current axioms.  Failure of one trial action or constitutive model can rule
out that realization within its tested scope; by itself it does not exclude
all other still-unspecified \ElEmEr{} dynamics.

\begin{itemize}
  \item \textbf{Closure.}  A proposed native local-$c$ law intended to explain
  massive matter must actually admit localized recurrent organizations with
  sufficient stability and isolation.  Failure under a specified law and
  admissible parameter range rejects that proposed matter realization;
  broader exclusion requires a genuinely exhaustive result over an
  independently defined admissible law class.

  \item \textbf{Relativity.}  Clocks, rods, signals, and interactions must
  recover the experimentally constrained operational Lorentz sector, without
  an observable leak of primitive succession or preferred-frame effects
  inconsistent with those observations.  Preferred-frame tests constrain specified
effective parameters, including those in the PPN description
\cite{Will2014Confrontation}; they are not direct measurements of $s$.
Local $c$ alone does not establish
  that recovery, and a model's failure is not a general theorem against every
  possible native realization.

  \item \textbf{Interaction.}  An explicit source, activation, and
  constitutive stress/current law must realize local field reorganization
  with a complete energy--momentum balance.  Merely writing a generic
  response template is not a derived microscopic force law.

  \item \textbf{Scale.}  The current framework treatment permits one common
  calibrated absolute scale and otherwise requires dimensionless
  organization.  A demonstrated need for unrelated, independently fitted
  sector scales would invalidate that \emph{one-scale strategy} for the
  proposed realization; it would not, by itself, disprove the primitive
  local-$c$ premise.

  \item \textbf{Explanatory independence.}  Agreement obtained only after
  importing the target theory's action, representation, or fitted coupling
  is a conditional consistency check, not a new \cmotion{} prediction.
  Genuine explanatory content requires additional constraints, parameter
  reduction, independently fixed selection rules, or a quantitative
  observable that could disagree with measurement.
\end{itemize}

A sharper empirical falsifier becomes available when a native law, its
admissible domain, and its free inputs are specified \emph{before} comparison
with the target outcome, so that it forces a quantitative prediction or
exclusion.  Failed predictions must not be repaired by unrestricted
post-hoc changes to the interaction or particle-specific calibration.
Conversely, neither success inside a borrowed effective model nor a negative
result for one candidate interaction constitutes a universal validation or
rejection of the unfinished framework.

The open microscopic requirements remain the \ElEmEr{} action and stress
functional; interaction source, activation and coupling laws; closure capture,
mode conversion, stabilization and controlled leakage; absolute scale origin;
operational Lorentz recovery without an observable $s$-leak; self-consistent
moving closures; and the deeper organization behind the observed particle
spectrum.  They remain explicit completion targets, not achievements inferred
from this list.

\section{Conclusion}
\label{sec:conclusion}

\cmotion{} starts from physical three-space, a locally universal propagation speed, and a
primitive before--after succession of physical states.  \ElEmEr{} is the underlying
field/state structure in which that propagation occurs.  Open propagation carries
change onward; recurrence provides physical cycles; sufficiently stable recurrence may
produce localized persistent organizations.  Ordinary time is then read from physical
clocks and networks of signals rather than inserted as a material fourth spatial-like
dimension.

The same field picture leads naturally to local interaction.  Incoming histories cease
to be primitive labels once they contribute to one complete local state.  Reorganization
of that state can redirect momentum flow, alter internal or modal structure, and in the
right circumstances help build recurrent return.  A newly produced change is proposed
to propagate onward causally at local $c$, so the framework does not require
instantaneous nonlocal updating of the field.

The proposed unity is physical, not merely terminological.  Open propagation
carries disturbance; activated interaction changes its local organization;
recurrent capture could sustain particle-like matter; cycles of matter can
serve as clocks; and the energy--momentum of those complete organizations
would govern their effective motion.  These are linked stages of one candidate
field dynamics, not separate fundamental ingredients or already proven stages.

Closure is the main constructive bet.  If local-$c$ motion can bend into a stable finite
return pattern, persistence need not be primitive.  If a self-consistent isolated pattern has a well-defined effective total
rest-frame energy, that quantity provides the conditional rest-mass interpretation;
reorganization of the whole return geometry under collective motion remains a
candidate physical route to inertia, not a native boosted-solution calculation.  The energy/action conjugacy and stress-steering relations show how these ideas
fit into a more precise dynamics once the underlying action is known.

The open problems are correspondingly specific.  The microscopic interaction that
captures and sustains closures is not yet derived; absolute scale still needs either a
single calibrated input or a dynamical origin; a globally aligned $s$, if it exists, must
remain operationally hidden; and the observed three-family structure has no present
mechanism.  Many of these gaps point toward the same upstream task: a physically
selected microscopic \ElEmEr{} law whose solutions can be compared across
sectors.  That shared dependency makes the law a potentially powerful route
to further progress, not an established mechanism or a guarantee that one
action automatically yields every missing physical result.  The gaps do not
erase the physical content already obtained: they identify exactly where the
proposed picture must acquire additional dynamics.

The framework can now be advanced by asking concrete questions rather than adding broad
new principles.  Can a local field law sustain finite recurrent closures?  Can a moving
closure reproduce Lorentz clock and rod behaviour without an $s$-leak?  Can one
interaction architecture generate the required sector-specific dynamics while keeping
energy--momentum accounting complete?  Can one absolute scale support the full physical
construction?  A fixed, independently specified microscopic law with discriminating
predictions would make these requirements decisive empirical tests; until then,
they are necessary completion criteria rather than a completed falsification theorem.

\newpage


\appendix

\section{Compact mathematical notes}
\label{app:math}

\subsection{Complete-state reversal and causal reachability}
\label{subsec:complete-state-reversal}
The following is the detailed reversal-compatible test behind the physical
source/restoration distinction in \cref{subsec:causal_reachability}.  It is a
candidate constructibility diagnostic, not a native arrow-of-time theorem.

The same picture raises a sharper question about physical reversal.  Suppose a localized
event generates a field/state change that propagates away along primitive local-$c$
characteristics.  A later change of the original source can generate further causal
changes, but those later changes also propagate locally at $c$.  Along the same
outward characteristic, a later disturbance cannot catch an earlier one if
both propagate at the same primitive speed.  More generally, under unchanged
admissible propagation conditions, a later emission from the same source cannot
precede the earlier emission's earliest causal arrival at a fixed destination;
interception along different selected paths is a distinct, geometry-dependent
question.  Reversing the motion or internal state of the visible source is not,
by itself, a reversal of the causal consequences that have already left it.

Here \emph{reversal-compatible} refers to a possible property of successive complete
physical states, not to time as a fourth physical dimension.  Schematically, if a
well-defined invertible state update $\mathcal U$ and an involution $R$ on the
complete physical state existed, the mathematical test would be
\begin{equation}
R^2=I,\qquad R\mathcal U R=\mathcal U^{-1}.
\end{equation}
The proposed $R$ would transform the appropriate propagation directions and
momentum-like data while retaining the full correlated field configuration.
This is only a reversal-compatibility \emph{test}: neither a native $\mathcal U$
nor its physical reversal map has been derived, and the notation does not
assume one globally aligned succession parameter.

An exact movie-reversal of such a history would require more than reversed source motion.
The already outgoing field structure would have to be replaced by the corresponding
precisely coordinated incoming structure, with the required directions, phases, internal
state, and correlations so that it converges on the appropriate later events.  If such an
incoming field is produced by other earlier sources, it has its own causal history and is
a new physical preparation; it is not the original source subsequently reversing the
field change that it previously generated.  In an otherwise closed system with no such
external preparation, the physical origin of the required incoming data is therefore
nontrivial.

This motivates a \emph{candidate causal-reachability constraint}.  A formal
reversal-conjugate solution and a later source-generated restoration are different
constructions.  When an explanation relies on precisely coordinated incoming field data,
those data should not be credited as a natively generated reversal mechanism merely
because they solve a mathematical field equation; their physical preparation remains a
separate question.  This is a constructibility guard, \emph{not} an established
requirement that every admissible instantaneous state have an independently certified
causal ancestry.  Whether the eventual native dynamics imposes a stronger restriction
on physical histories remains open.

The proposal is especially natural in the present primitive ontology.  Physical changes
propagate through physical three-space; they are not treated as a material four-dimensional
spacetime history that can simply be geometrically ``rewound'' by reversing a local source.
Once a causal change has propagated away at local $c$, any later influence from the same
interior source remains subject to the same propagation limit.

This is presently a candidate constraint, not a derived irreversibility theorem.
Local-$c$ transport alone does not exclude a complete reversal-conjugate field state;
the decisive missing test is whether the eventual native generation, interaction, and
absorption laws permit the fully reversed process.  A stronger physical-history rule
might emerge from those laws or additional admissibility conditions, but is not assumed
here.  In particular, this proposal does not derive entropy increase or the
thermodynamic arrow.

\subsection{Conditional rest mass and the complete stress-trace test}
\label{subsec:stress-trace-mass}
The standard Laue/virial constraint \cite{Laue1911Dynamik,Giulini2018LaueRevisited}
uses a complete total stress tensor, not the
stress of an isolated carrier component and not a new \ElEmEr{} stress law.

A further \emph{Conditional} stress-trace restriction makes the missing microscopic
binding problem more specific.  Suppose a complete, isolated, localized rest
configuration possesses a symmetric conserved total stress--energy tensor
$T^{\mu\nu}$, sufficiently decaying boundary terms, and a stationary Laue
condition, or a bounded recurrent motion with a valid vanishing virial-average
boundary term.  In signature $(+---)$, the integrated spatial stresses then
vanish (or vanish in the stated cycle average), giving
\begin{equation}
 \boxed{Mc^2=E_{\rm rest}
   =\int\!\langle T^{00}\rangle\,\dd^3x
   =\int\!\langle T^\mu{}_{\mu}\rangle\,\dd^3x.}
 \label{eq:conditional_stress_trace_mass}
\end{equation}
Here $\langle\cdot\rangle$ denotes the appropriate stationary value or
bounded recurrence/virial average, not a thermodynamic ensemble average.
The \emph{local} energy density remains $T^{00}$, not the trace.  Thus a
fully traceless \emph{complete} stress tensor cannot, under these assumptions,
support a positive rest energy; nonzero integrated trace is a necessary condition
in this class, not a proof of capture or a unique binding mechanism.  Boundary-
supported, radiating, nonisolated, or incompletely accounted interaction-stress
systems require separate analysis.  No native $T^{\mu\nu}$ or physical closure
mass has been derived by this standard conservation/virial identity.

\subsection{Pure-ray transparency versus mixed-state Lorentz compatibility}
\label{subsec:no-s-leak-countercontrol}
This algebraic countercontrol tests a restricted two-channel interaction;
it does not constitute a native dynamical or moving-clock model.

A restricted algebraic
countercontrol illustrates the gap: for opposed channel intensities
$A=|R|^2$, $B=|L|^2$, the interaction
$V_{\rm leak}=\lambda AB(A+B)$ and its first variation vanish on either pure-ray
axis, yet reciprocal reweighting $A\mapsto e^\eta A$,
$B\mapsto e^{-\eta}B$ changes $V_{\rm leak}$ for mixed states.  Thus even exact
open-ray transparency does not select a boost-compatible interaction.  This is
an illustrative model-class counterexample, not a proposed native \ElEmEr{}
action or evidence for an observed preferred frame.

\subsection{A complementary charged-lepton conversion test}
\label{subsec:conversion-cycle}
The following conditional example concerns one proposed exactly conserved
count, not the general charge or energy bookkeeping of field theory.

\paragraph{Complementary electron--muon conversion-cycle illustration (conditional).}
A different stress test of permanent parcel bookkeeping comes from an
energetically accessible charged-lepton conversion sequence:
\begin{align}
e^-+e^+ &\longrightarrow \mu^-+\mu^+,\\
\mu^- &\longrightarrow e^-+\bar\nu_e+\nu_\mu,\\
\mu^+ &\longrightarrow e^++\nu_e+\bar\nu_\mu.
\end{align}
The net outgoing asymptotic pattern is an electron--positron pair and four
neutrinos, rather than just the original two charged leptons.  The conversion
requires sufficient center-of-momentum collision energy; repeating the sequence
requires fresh energy input to produce another muon pair.  It is therefore
\emph{not} an energy-producing infinite loop.  If a candidate parcel ontology
assigns the same invariant nonzero count to every outgoing charged lepton and
neutrino, and its external energy supply carries no such count, the change
from two counted particles to six contradicts that conservation premise.  But the
example alone does not exclude species-dependent or signed charges, a
count-carrying preparation field, or a general continuous conserved field
functional.  It is a \emph{conditional conversion/branching illustration}, not a
second unconditional no-go theorem.

\subsection{Constitutive selection beyond consistency checks}
\label{subsec:selector-controls}
The reduced selector models delimit what follows from shared consistency
properties and endpoint conditions; they do not select a complete native law.

The native-dynamics selector research gives a sharper reason that compatibility
alone is insufficient.  Within an explicit reduced smooth control, two laws
share the tested local-$c$ kinematics, free and recurrent endpoints, conserved
budgets, and other consistency guards yet differ in the physical transition
history.  This demonstrates a \emph{scoped} constitutive underdetermination,
not a no-go theorem against selecting any native law.  A future completion must
supply (i) an active source or activation condition from the complete physical
state, (ii) the selected physical response direction, (iii) its rate/strength
normalization in primitive succession, and (iv) branch selection or weights if
successors are not unique.  An action is one possible representation of such
a law, not a necessary assumption about its fundamental form.  A positive
reduced selector control becomes rigid only after independent source-composition,
response-channel, and normalization inputs are stipulated; none of those inputs
is thereby derived for native \ElEmEr{} dynamics.

\subsection{Scope of local geometry and null-flow exchange}
\label{subsec:null-flow-proof}
The three-space arclength $\dd\ell^2=h_{ij}\dd x^i\dd x^j$, the local
orthonormal three-vector formulas, and any derivative $\partial_s$ used on a spatial
neighbourhood require spatial and local-evolution chart assumptions.  The
simple displayed balance equations additionally use a chart with no omitted
metric/volume evolution terms; otherwise those terms must enter the balance.
These assumptions do not
follow from the isolated characteristic relation $\dd\ell=c\,\dd s$ and do not
prove a globally synchronized succession parameter.

For one resolved null flow, write
$\mathbf p_{\rm car}=\varepsilon\mathbf u/c^2$ and
$\boldsymbol\Pi_{\rm car}=\varepsilon\mathbf u\otimes\mathbf u/c^2$.
Expanding Eq.~\eqref{eq:resolved_null_momentum_balance} and using
Eq.~\eqref{eq:resolved_null_energy_balance} gives
\begin{equation}
 \frac{\varepsilon}{c^2}D_s\mathbf u
 +\frac{q}{c^2}\mathbf u=\mathbf f.
\end{equation}
Projection along and perpendicular to $\mathbf u$ produces respectively
$q=\mathbf u\cdot\mathbf f$ and Eq.~\eqref{eq:force_balance}.  Writing
$\mathbf f=-\nabla\cdot\boldsymbol\Sigma_{\rm int}$ additionally requires
control of the interaction-sector momentum-density evolution.  Nothing here
proves that a general interacting ElEmEr field is a single null stream.

\subsection{Why the projected law preserves speed}
\label{subsec:speed-proof}
For Eq.~\eqref{eq:projected_steering},
\begin{align}
\frac{1}{2}\frac{\dd}{\dd s}|\mathbf u|^2
&=\mathbf u\cdot\frac{\dd\mathbf u}{\dd s}\\
&=\mathbf u\cdot\mathbf P_\perp\mathbf g_{\rm st}
+\mathbf u\cdot\left(\frac{\mathbf u}{c}\times\mathbf h_{\rm st}\right)=0.
\end{align}
Thus $|\mathbf u|=c$ is preserved within the stated representation.

\subsection{Projective all-null theorem}
\label{subsec:all-null-proof}
Decompose the linear generator $A=F+S$ into its $g$-skew-adjoint and $g$-self-adjoint
parts: $g(U,FV)=-g(FU,V)$ and $g(U,SV)=g(SU,V)$.  For every null $K$,
$g(K,FK)=0$, so the infinitesimal cone-tangency hypothesis $g(K,AK)=0$
reduces to $g(K,SK)=0$.  To verify the dimension range, choose an orthonormal
basis with $g=\operatorname{diag}(1,-1,\ldots,-1)$ and write
$B_{\mu\nu}=g_{\mu\alpha}S^\alpha{}_{\nu}$, a symmetric matrix.
For each spatial basis vector $e_i$, the null vectors $(1,\pm e_i)$ give
$B_{00}+B_{ii}\pm2B_{0i}=0$.  Hence $B_{0i}=0$ and
$B_{ii}=-B_{00}$.  If there are at least two spatial dimensions, using
$(1,(e_i+e_j)/\sqrt2)$ also gives $B_{ij}=0$ for $i\ne j$.
With just one spatial dimension there are no such off-diagonal spatial entries,
so the first step already completes the proof.  Thus $B=\phi g$ and
$S=\phi I$ in every Lorentzian dimension at least two, giving $A=F+\phi I$.
For example, in $1+1$ dimensions a boost generator sends the two null vectors
$(1,\pm1)$ to opposite multiples of themselves: normalization can change even
when projective direction cannot.  The
hypothesis constrains an infinitesimal generator; it does not require $AK$ to be
null.  The result is conditional on one common real local-linear operator acting
on \emph{all} null directions and does not identify a unique primitive force field.

\subsection{All-direction sphere test for one steering two-form}
\label{subsec:sphere_response_test}
At a fixed local state and with a physically fixed null-tangent normalization,
let $\mathbf a_\perp(\mathbf n)$ be the measured or computed \emph{transverse}
response for each unit direction $\mathbf n$, holding the independent internal
state data fixed.  Within the common real local-linear/projective class of
Eq.~\eqref{eq:all_null_A}, a universal six-component response requires
\begin{equation}
 \mathbf a_\perp(\mathbf n)
   =(I-\mathbf n\mathbf n^{\mathsf T})\mathbf g_{\rm st}
   +\mathbf n\times\mathbf h_{\rm st}\quad (|\mathbf n|=1).
 \label{eq:sphere_response_form}
\end{equation}
With normalized spherical averaging and
$\langle n_i n_j\rangle_{S^2}=\delta_{ij}/3$, the only possible coefficients are
\begin{equation}
 \mathbf g_{\rm st}=\frac32\langle\mathbf a_\perp\rangle_{S^2},\qquad
 \mathbf h_{\rm st}=-\frac32\langle\mathbf n\times\mathbf a_\perp\rangle_{S^2}.
 \label{eq:sphere_response_reconstruction}
\end{equation}
The residual
\begin{equation}
 \mathbf R(\mathbf n)=\mathbf a_\perp(\mathbf n)
 -(I-\mathbf n\mathbf n^{\mathsf T})\mathbf g_{\rm st}
 -\mathbf n\times\mathbf h_{\rm st}
 \label{eq:sphere_response_residual}
\end{equation}
must vanish for \emph{all} directions for one universal angular two-form
representation in that class.  Failure rejects this representation of the
tested response; it does not by itself reject all nonlinear or internally
state-dependent steering laws.  The test neither derives Maxwell dynamics nor
removes the complete exchange-current and momentum-density assumptions of
Eq.~\eqref{eq:resolved_null_momentum_balance}.

\subsection{Energy--action conjugacy versus product law}
\label{subsec:conjugacy-proof}
For the admissible smooth family $z(a)$ of constrained stationary states in
\cref{subsec:energy_action_layers}, stationarity at each member means
$\delta E[z(a)]=J_0\Omega(a)\delta Q[z(a)]$ on variations preserving the
remaining charges and boundary conditions.  Apply this equality to the tangent
$z'(a)$ to obtain $E'(a)=J_0\Omega(a)Q'(a)$, and hence
$\dd E/\dd Q=J_0\Omega$ wherever $Q$ is a valid local family coordinate.
One does not differentiate $\Omega$ in the fixed-multiplier field variation.
This is the stationary-family construction used in Q-ball analysis
\cite{Coleman1985Qballs}, not an existence or stability proof for native closures.  In this variational statement,
$E$ must generate translations of the same specified evolution parameter whose phase
rate is $\Omega$; neither need already have a native observer-clock interpretation.
Integrating along the family yields Eq.~\eqref{eq:energy_integral}.
If a product identity $E(Q)=J_0\Omega(Q)Q$ were to hold for all $Q$ in a nonzero
smooth charge interval, its derivative would give
$\dd E/\dd Q=J_0(\Omega+Q\,\dd\Omega/\dd Q)$, hence the conjugacy relation requires
$Q\,\dd\Omega/\dd Q=0$.  Thus $\Omega$ must be constant on that interval and the
energy zero chosen consistently.  A special equality at one value of $Q$ is not
a universal nonlinear identity.

\section{Claim-status register}
\label{app:status}

This appendix records the scientific status of the claims summarized in this paper.  It is self-contained: each entry states the status and the scope required for that status, including imported or companion-paper dependencies where relevant.

The status vocabulary follows the programme's Derivation Ledger:
\emph{Assumed} denotes a provisional active premise; \emph{Conditional}, a
conclusion requiring an explicit premise or model class not itself derived;
\emph{Structural}, an exact or established organizational/accounting statement
that does not supply microscopic dynamics; \emph{Conjectural}, a motivated but
underived mechanism or interpretation; \emph{Negative}, an exclusion within its
stated scope; and \emph{Open}, an unresolved question.  \emph{Open (sharp)}
is an isolated missing bridge, not greater confidence.  \emph{Deferred} marks
an item outside the present publication critical path while retaining its
scientific status.  A slash preserves distinct aspects or dependencies rather
than creating a new status.  These classifications do not promote imported
mathematics to a native dynamical derivation.

\small
\begin{longtable}{@{}P{4.9cm}P{3.0cm}P{6.8cm}@{}}
\toprule
Claim & Status & Scope / limitations \\
\midrule
\endfirsthead
\toprule
Claim & Status & Scope / limitations \\
\midrule
\endhead
Physical three-space; time not primitive spatial-like dimension & \statusAssumed{} & Local positive spatial metric/orthonormal-chart conventions are required for the displayed vector formulas; neither a unique global geometry nor a fourth material time dimension is derived. \\
Primitive local-$c$ propagation & \statusAssumed{} & Kinematic premise; not an energy amount or Lorentz proof. \\
Primitive succession $s$ with $\dd\ell=c\,\dd s$ & \statusAssumed{} / \statusStructural{} & Local evolution/order measure; PDEs additionally assume compatible local evolution charts. \\
One globally aligned $s$ & \statusSharpOpen{} & Candidate foundational possibility; not derived or operationally readable. \\
Recurrent clocks and operational time & \statusStructural{} / \statusConditional{} & Clock readings are physical process counts; moving-clock dynamics remain conditional. \\
Operational Lorentz relations & \statusConditional{} & Require no observable $s$-leak and correct closure/rod/synchronization dynamics. \\
Field-first complete-local-state interaction ontology & \statusStructural{} & Ancestry labels are bookkeeping, not primitive persistent parcels. \\
Candidate Interaction--Response architecture & \statusConjectural{} / \statusStructural{} & Opportunity is not sufficient activation, attraction, binding, or mode conversion. \\
Candidate causal propagation clause & \statusConjectural{} & Generated changes propagate successor-directed at local $c$; generation law remains open. \\
Candidate causal-reachability constraint
& \statusConjectural{} / \statusOpen{}
& Formal reversal means a hypothetical complete-state succession symmetry, not a fourth time dimension. Later source-only restoration is not complete field reversal; native generation/absorption symmetry remains open. \\
Common response architecture across sectors & \statusConjectural{} / \statusStructural{} & Shared template only; no universal force is derived. \\
Projected speed-preserving steering law & \statusConditional{} & Kinematic representation; source/action/constitutive law open. \\
Stress-divergence steering & \statusConditional{} & Resolved rank-one null-flow and exchange-balance assumptions; stress-only form needs interaction momentum-density control. Not a generic coherent-field law. \\
All-null $A=F+\phi I$ theorem & \statusConditional{} & One real common infinitesimal operator, all null directions, Lorentzian dimension at least two; $F$ can rescale as well as reorient. No native physical two-form or amplitude law derived. \\
Universal angular sphere-response diagnostic & \statusConditional{} & Six-component local-linear response test after normalization and full directional sampling; not a native stress law or Maxwell derivation. \\
Energy/action separation & \statusStructural{} / \statusConditional{} & Energy requires autonomous Hamiltonian/action; action charge requires exact phase/action symmetry. \\
Shared microscopic law as an upstream research dependency & \statusOpen{} / \statusDeferred{} & A selected native action or equivalent local law could constrain several sector problems; reduced selector controls leave source, response, rate and branch selection unresolved. No complete law or selection is derived. \\
$\dd E/\dd Q=J_0\Omega$ & \statusConditional{} & Smooth recurrent symmetry family with fixed other charges/topology; $E$ and $\Omega$ use one specified evolution generator, not an automatically observer-clock frequency. \\
Generic $E=J_0\Omega Q$ & \statusNegative{} & Not a general nonlinear-closure identity; an interval-wide product law requires constant $\Omega$ at fixed other charges and a suitable energy zero. \\
Generic finite-wave positive action current & \statusSharpOpen{} & Positive conserved current and justified decomposition for arbitrary finite coherent/interacting waves remain unresolved (\cref{subsec:energy_action_layers}); resolved-mode factorization is insufficient. Quantum normalization remains with Quantum Foundations. \\
Universal normalization slot $J_0$ & \statusStructural{} & Common positive rescaling of a complete classical action preserves its equations but rescales action-derived charges and stress; the final physical normalization remains unresolved. \\
Identification $J_0=\hbar$ & \statusAssumed{} / \statusSharpOpen{} & Treated in the Quantum Foundations companion paper; not derived here. \\
Open/closed/extended organization & \statusStructural{} & Overlapping, nonexhaustive descriptors of field organization, not an exclusive partition or a derivation of spin, charge, family, or binding. \\
Return constraints can discretize a specified closure problem & \statusStructural{} & Scalar phase is a model limit; full internal holonomy must be supplied, not presumed spinorial. No derived $E=hf$, empirical spectrum, or family count. \\
No identical nonzero additive count per outgoing photon & \statusNegative{} & Fixed initial quantity, identical constant photon weight, exact additivity/conservation, and no uncounted carriers imply $w_\gamma=0$ for differing multiplicities. QED predicts $3\gamma/5\gamma$; five-photon measurements alone are not decisive against zero. The $2\gamma/4\gamma$ comparison has stronger experimental support. Conversion-cycle illustration remains Conditional. \\
Mass as invariant proper-return energy & \statusConditional{} & $E_{\rm return}$ is the total rest-frame energy of a complete isolated recurrent solution, assuming conserved energy--momentum and a rest frame; $Mc^2=E_{\rm return}$ is not a native computed mass. \\
Inertia as full-return energy--momentum response & \statusConditional{} & Whole-pattern reorganization is an interpretation; the inertial coefficient currently follows only from an assumed effective Lorentz dispersion. Native boosted closure remains Open. \\
Rest/inertial mass equality & \statusConditional{} & Inherited from an effective Lorentz-covariant energy--momentum relation, not a derivation of gravitational equivalence or free fall. \\
Conditional Laue/virial stress-trace identity & \statusConditional{} & Standard identity for complete symmetric conserved isolated stress with vanishing boundary/virial terms; nonzero integrated trace necessary for positive rest energy in this class, not a native binder or mass prediction. \\
Pure-ray transparency does not ensure mixed-state Lorentz compatibility & \statusNegative{} & In the opposed-channel algebraic control, $\lambda AB(A+B)$ and its first variation vanish on pure axes but reciprocal reweighting changes mixed-state interaction (\cref{subsec:no-s-leak-countercontrol}); no general exclusion of native laws. \\
Weak-field gravity as elementary-closure binder & \statusNegative{} & Measured weak-field $G$, elementary-particle energies/radii and tiny $GE/(Rc^4)$; not a strong-gravity exclusion (\cref{subsec:binder}). \\
Direct same-charge Abelian contractile binder & \statusNegative{} & Standard positive-energy classical vector self-energy with same-charge repulsion and no additional cohesive stresses; not arbitrary nonlinear, scalar or coherence channels (\cref{subsec:binder}). \\
Native microscopic binder & \statusOpen{} / \statusDeferred{} & Scientific completion deferred; no solved universal binder claimed. \\
Absolute scale generation & \statusOpen{} / \statusDeferred{} & One calibrated energy-valued $\Lambda_C$ permitted as fallback; $J_0c/\Lambda_C$ is a conditional dimensional reference, not a derived closure size. \\
Exactly-three-family mechanism & \statusOpen{} / \statusDeferred{} & Explicit non-derivation; detailed exclusions are given in the relevant companion papers. \\
Global phase $\not\Rightarrow$ local gauge dynamics & \statusStructural{} / \statusConditional{} & Local gauge connection/action remains sector-specific. \\
\bottomrule
\end{longtable}
\normalsize

\section{Terminology glossary}
\label{app:glossary}

\begin{description}[style=nextline]
\item[Elemental Motion Resonance (ElEmEr)]
The proposed underlying dynamical field/state structure of the framework.  Open
propagating modes, recurrent closures, and extended organizations are different forms of
its organization.  Its final microscopic action remains to be derived.

\item[ElEmEr wave]
A travelling or standing wave-like excitation pattern of the underlying field, not
an individually conserved material parcel or independent substance.

\item[ElEmEr field excitation]
An organized disturbance of the underlying local field state; need not be a travelling
wave or an autonomous recurrence.  A field \emph{value} is a value of a chosen local
variable and need not specify the complete dynamical state.

\item[Open propagating mode]
An admissible local-$c$ propagating component that does not itself constitute an
autonomous recurrent closure; such a component may coexist with closed or extended
organization elsewhere in the complete field state.

\item[Closure]
A localized recurrent organization of \ElEmEr{} propagation.  Closure denotes the
organization itself; capture history, microscopic binder, stability, and empirical
particle identity are additional questions.

\item[Extended organization]
A causally prepared physical organization whose relevant state spans separated regions;
this does not by itself imply superluminal signalling.  It may encompass localized
recurrent closures and open components rather than being a disjoint third substance.

\item[Primitive succession $s$]
The assumed primitive before--after ordering of successive physical states, distinct from
physical clock readings and observer coordinate time.

\item[Interaction opportunity]
A local state condition in which the system has left a valid free-propagation class and
a response channel may become available.  Actual activation depends on the
sector-specific dynamics.

\item[Steering]
A local change of propagation direction within a speed-preserving response
representation.

\item[Binder]
The microscopic dynamical mechanism that would supply the capture and stability needed
for an autonomous closure.

\item[$\Lambda_C$]
The one permitted framework-wide absolute-energy-scale fallback in the present
work-in-progress formulation.  The preferred long-term outcome is a dynamically
generated scale.  Dimensional reference lengths additionally require an action
normalization such as $J_0$; they are not predicted closure radii.
\end{description}

\newpage

\printbibliography


\end{document}
